\documentclass[a4paper,10pt]{article}
\usepackage[utf8]{inputenc}
\usepackage{graphicx}
\usepackage{subfigure}
\usepackage{fancyhdr}
\usepackage{listings}
\usepackage{ gensymb }


%opening

\pagestyle{fancy} %eigener Seitenstil
\fancyhf{} %alle Kopf- und Fußzeilenfelder bereinigen
\fancyhead[L]{SEP Abgabegespräch SS2015/16} %Kopfzeile links
\fancyhead[C]{} %zentrierte Kopfzeile
\fancyhead[R]{\includegraphics[width=1cm]{../../template/tugraz_bw}} %Kopfzeile rechts
\renewcommand{\headrulewidth}{0.4pt} %obere Trennlinie
\fancyfoot[C]{} %Seitennummer
\renewcommand{\footrulewidth}{0.4pt} %untere Trennlinie


\lstset{language=C++, basicstyle=\ttfamily\tiny}  

\begin{document}

\section*{Notenschnitt-Rechner}

Schreibe ein einfaches Programm um den Notenschnitt mehrerer Studenten zu berechnen. Die Studenten werden mit Namen erstellt, und jede Lehrveranstaltung besitzt einen Namen und die Endnote.


\begin{description}
 \item[create $<$studentName$>$] Erstellt einen neuen Studenten mit dem Namen $studentName$. 
 \item[add $<$studentName$>$ $<$subjectName$>$ $<$subjectGrade$>$] Der Student mit Namen $studentName$ bekommt die Note $<$subjectGrade$>$ für die Lehrveranstaltung $<$subjectName$>$.
 \item[show] Listet alle Studenten mit deren Noten auf. 
 \item[average] Listet den Notendurchschnitt aller Studenten auf.
\end{description}
Bei Fehlerfällen (Falsche Parameter, kein Held existiert...) soll nichts ausgegeben werden.
Bei der Eingabe von \texttt{quit} oder EOF soll das Programm beendet werden.

\subsection*{Beispiel}
\begin{verbatim}
> create SuperStudent
> create HauptsacheDurch
> add SuperStudent Analysis 1.0
> add SuperStudent SEP 2.0
> add HauptsacheDurch NRLA 4.0
> add HauptsacheDurch Messtechnik 5.0
> show
HauptsacheDurch
NRLA: 4
Messtechnik: 5
SuperStudent
Analysis: 1
SEP: 2
> average
Grade average of HauptsacheDurch: 4.5
Grade average of SuperStudent: 1.5
> quit
\end{verbatim}

\subsection*{Aufruf}
\texttt{./grade-average}

\subsection*{Bewertung}
Dokumentation \& Programmierstil,  Dokumentation \& Programmierstil, und Struktur \& Korrektheit (je 1 Punkt)
\newpage


\end{document}
